This work studied the BCD-MMSE algorithm in the uplink of a massive MIMO system with decentralized processing and colored noise. The focus was on the inversion of the matrix $\mathbf G_c$ at each DU, performed in fixed-point arithmetic on an FPGA. \textcolor{red}{Essa inversão ainda não tinha implementação em hardware (Capítulo~\ref{chap:relatedwork}), e este trabalho apresenta uma.} A Python-based simulator was used to assess the impact of the numerical format on the SER curves. Subsequently, the inversion module was implemented in hardware, with one architecture for each test condition.

% FONTE: Tabela tab:formatos-escolhidos e Seções sec:resultados-formato, subsec:perfis-mobilidade,
%   sec:resultados-curvas (formatos.tex). Mobilidade: formatos.tex, "little impact on the SNR
%   increase"; só com baixa mobilidade, 2 bits a menos em dois casos (G2 3 dB, G3 0,1 dB).
The format depends on the group and on the SNR increase limit. With all configurations, the 3~dB limit requires $\langle 24,12\rangle$ in groups 1 and 3, and $\langle 28,14\rangle$ in group 2, while the 0.1~dB limit requires $\langle 26,12\rangle$ in group 1 and $\langle 30,14\rangle$ in groups 2 and 3 (Table~\ref{tab:formatos-escolhidos}). \textcolor{blue}{The formats have 2 to 6 fewer bits without the configurations with an IoTh of 20~dB.} QPSK is the modulation most sensitive to the format, as it operates at a lower SNR, where the elements of $\mathbf G_c$ overflow the integer part (Section~\ref{sec:resultados-curvas}). With 16-QAM, the effect of the format is almost negligible, and the difference from LMMSE comes almost entirely from decentralization (Section~\ref{sec:resultados-formato}). \textcolor{blue}{The mobility profile has little effect on the SNR increase for the same format.} \textcolor{blue}{The low-mobility conditions alone would allow a format 2 bits smaller in only two cases, so the high-mobility profile defines the format (Section~\ref{subsec:perfis-mobilidade}).}

% FONTE: hardware.tex (tab:recursos-faixa, tab:potencia-dinamica, fig:energia-inversao,
%   subsec:resultados-etapas).
All architectures \textcolor{red}{that fit in the FPGA} meet the cycle limit of their test condition at 100~MHz (Section~\ref{subsec:resultados-recursos}). Resource usage and dynamic power follow the parallelism. The dynamic power stays below 1.1~W with loop pipelining, dataflow, and matrix batching, and reaches up to 4.7~W with systolic arrays (Section~\ref{subsec:resultados-potencia}). The energy per inversion, however, is higher in the \textcolor{red}{optimized} architectures with less parallelism, those of TDLA30-5 and TDLA30-10. \textcolor{blue}{The LU decomposition and the forward substitution, which contain the dividers, account for most of the resources and power in all architectures (Section~\ref{subsec:resultados-etapas}).} \textcolor{red}{In $32 \times 32$ (group 2), only the version without optimizations was implemented, and its interval between inversions is above all the cycle limits (Section~\ref{subsec:resultados-recursos}).}

% FONTE: compromisso.tex (tab:escolha-condicao, fig:quadrantes-g1/g3). TDLB100-400 sem <30,14>
%   (hardware.tex); <28,14> fica em 0,13 dB nessa condição.
In the trade-off with the SNR increase, the format changes power and energy less than the architecture does (Section~\ref{subsec:resultados-compromisso}). \textcolor{blue}{A format smaller than the one of the group is sometimes sufficient, since this analysis is done for each test condition and each IoTh.} The exception is TDLB100-400 in group 3, with an IoTh of 20~dB: the $\langle 30,14\rangle$ format does not fit in the FPGA, and no smaller format stays within 0.1~dB.

% FONTE: latencia-vazao.tex (tab:atraso-bloco, ITU 1 ms, outras tarefas da BS).
The architectures were sized by throughput rather than by delay. Under low-mobility conditions, they distribute the inversions over the coherence time and use less area and power, but the last inversion of a block finishes later. The collection of the $N$ noise samples, required by the estimation method, also adds delay. \textcolor{blue}{Only TDLA30-300 and TDLB100-400 stay below the 1~ms latency expected by the ITU-R with data collection and inversion (Section~\ref{sec:resultados-atraso}), without counting the other BS tasks.}

\textcolor{blue}{Thus, the fixed-point inversion of $\mathbf G_c$ in $16 \times 16$ fits in the XCU50 and meets the throughput of all test conditions, with formats within 3~dB and within 0.1~dB of the floating-point SER. The only exception is the 0.1~dB limit in TDLB100-400 of group 3, with an IoTh of 20~dB.}

